\noindent\hspace*{\parindent}Zu den einzelnen Positionen der vorstehenden Bilanzrechnung gebe ich die folgenden näheren Erläuterungen.

Zu (1). Der Kapitalwert des festen Teils der Witwenpensionen für sämtliche 15 jetzt lebenden Witwen (zu 200 Thaler für jede) ist die Summe der jetzigen Kapitalwerte dieser Pensionen, für jede einzelne Witwe nach ihrem Lebensalter berechnet, unter Berücksichtigung möglicher Waisenpensionen, wo jetzt minderjährige Kinder vorhanden sind. Ich führe diese Werte einzeln auf: Die Witwen sind nach der vollständigen Reihenfolge seit der Gründung der Anstalt nummeriert.
\begin{center}
\footnotesize\setlength{\tabcolsep}{3pt}
\begin{tabular}{r|l|r l||r|l|r l||r|l|r l}
31&\textsc{Sch.}&1066&Thaler&55&H.&1834&Thaler&65&H.&2003&Thaler\\
44&F.&1622&&58&G.&2730&&66&H.&1172&\\
50&H.&2849&&59&W.&1862&&68&M.&2923&\\
52&P.&1479&&62&B.&1003&&69&D.&2897&\\
53&M.&1812&&63&G.&1309&&70&L.&2224&
\end{tabular}
\end{center}
\noindent\hspace*{\parindent}Die Zahlen (2) und (7) sind mit folgendem Verfahren ermittelt worden, dessen Begründung in der Denkschrift von 1845 zu finden ist. Unter den 53 Mitgliedern, die gegenwärtig die Genossenschaft bilden, sind derzeit 47 verheiratet. Für jedes von ihnen ist nach dem Alter des Mannes und der Frau der jetzige Kapitalwert sowohl der vom Mann zu leistenden Beiträge (zu 15 Thaler jährlich) als auch des festen Teils der Witwenpension berechnet worden, die der Frau zusteht, falls sie den Mann überlebt (zu 200 Thaler). So verstanden ergibt sich die Summe der Beiträge zu 7947 Thaler, die Summe der Pensionen zu 45364 Thaler. Um die 6 jetzt unverheirateten Mitglieder mit zu berücksichtigen, werden diese Zahlen mit \(\tfrac{53}{47}\) multipliziert, woraus die Beiträge \(=8961\) Thaler und die Pensionen \(=51155\) Thaler hervorgehen: Die erste Zahl ist die obige Position (7). Die zweite ergibt, wegen der Waisenpension um \(\tfrac{1}{7}\) erhöht, 58463 Thaler, die Position (2).

Die Begründung dafür, bei den Waisenpensionen den Bruch \(\tfrac{1}{7}\) anstelle des 1845 verwendeten Bruchs \(\tfrac{1}{6}\) anzunehmen, wird bei der Erläuterung der Positionen (3) und (8) gegeben werden.

Ich führe noch die Ergebnisse der obigen Rechnung für die einzelnen 47 verheirateten Mitglieder auf. Die Nummerierung folgt der Reihenfolge des Eintritts in die Anstalt; die Zahlen der dritten Spalte sind die Beiträge, die der vierten die Pensionen.

\SourceLeaf{171}
\begin{center}
\footnotesize\setlength{\tabcolsep}{2pt}
\begin{tabular}{r|l|r|r||r|l|r|r||r|l|r|r}
135&O.&102 Thl.&1179 Thl.&170&R.&124 Thl.&1261 Thl.&199&F.&167 Thl.&937 Thl.\\
138&C.&82&1402&171&C.&132&955&200&L.&211&748\\
139&U.&137&730&173&F.&168&931&201&W.&194&749\\
140&H.&134&1199&174&H.&163&1354&202&M.&201&991\\
141&L.&113&833&175&R.&195&996&203&B.&167&751\\
147&R.&133&1188&176&L.&190&1122&204&E.&201&697\\
149&G.&127&1439&177&W.&176&936&205&H.&195&790\\
151&O.&173&870&180&T.&177&887&206&H.&188&854\\
153&B.&154&1271&185&W.&198&919&207&W.&200&783\\
155&K.&145&727&186&H.&169&1134&208&E.&175&1278\\
156&v. S.&146&662&187&W.&180&748&209&B.&154&1051\\
160&H.&161&1173&189&H.&172&953&211&S.&195&1164\\
161&B.&156&682&192&R.&188&782&212&T.&181&895\\
164&W.&159&1155&193&G.&206&896&213&W.&200&811\\
167&S.&190&835&195&D.&191&822&214&B.&188&915\\
169&Z.&180&943&196&B.&207&968&&&&
\end{tabular}
\end{center}
